\documentclass[11pt]{article}
% Uses the canonical latex-docs semantic style in-repository. The local
% lightweight fallback is for independent compilation of this distribution.
\IfFileExists{technical-implementation.sty}{\usepackage{technical-implementation}}{\usepackage{../support/aaoa-standalone}}
\usepackage{mathtools}
\setDocTitle{Chapter-to-Operations Crosswalk}
\setDocSubtitle{Nine-chapter map and artifact traceability}
\setDocVersion{1.0}
\setDocStatus{Educational reference / unapproved working material}
\setDocOwner{Algorithm Engineering}
\setDocAuthor{Original companion material}
\setDocSummary{Exact printed chapter/topic map plus operational work products and prerequisite dependencies.}
\providecommand{\Expect}{\mathbb{E}}
\providecommand{\Prob}{\mathbb{P}}
\providecommand{\Var}{\operatorname{Var}}
\providecommand{\bigoh}{\mathcal{O}}
\begin{document}
\makedoctocpage

\section{Reading crosswalk and source scope}
The sections below correspond to the chapter and section taxonomy of Sedgewick--Flajolet, second edition (2013). Page references are the book's \emph{printed pages}, not PDF viewer page indices. The lab IDs and assessment links are original suite mappings and are not part of the source text.
\begin{longtable}{p{0.07\linewidth}p{0.24\linewidth}p{0.26\linewidth}p{0.25\linewidth}}\toprule Ch. & Primary section sequence & Operational application & Deliverables \\ \midrule
1 & 1.1--1.8 (pp. 3--33): models, quicksort, randomness & Define cost/input model; derive quicksort predictions & SG01, WB01, LAB-PY01 \\
2 & 2.1--2.7 (pp. 43--80): first/higher-order and divide-and-conquer recurrences & Solve exact operation counts & SG02, WB01, LAB-CPP01 \\
3 & 3.1--3.11 (pp. 92--140): OGF/EGF/PGF, bivariate GFs & Coefficient extraction and probability models & SG03, WB02, LAB-PY02 \\
4 & 4.1--4.9 (pp. 153--211): asymptotics, Euler--Maclaurin, Laplace & Derive approximations and errors & SG04, WB03 \\
5 & 5.1--5.5 (pp. 220--247): symbolic methods and coefficient asymptotics & Translate classes into GFs and growth laws & SG05, WB02 \\
6 & 6.1--6.15 (pp. 258--331): trees, BSTs, path length, Lagrange inversion & Random BST and Catalan-tree comparison & SG06, WB04, LAB-PY03 \\
7 & 7.1--7.9 (pp. 347--406): permutations, inversions, cycles & Insertion-sort shift counts & SG07, WB04, LAB-CPP02 \\
8 & 8.1--8.9 (pp. 416--465): strings, automata, tries & Prefix-search cost models & SG08, LAB-PY04 \\
9 & 9.1--9.8 (pp. 474--532): occupancy, hashing, mappings & Bucket-length and probe experiments & SG09, WB04, LAB-CPP02 \\
\bottomrule\end{longtable}
\section{Method-before-structure dependency map}
\begin{align*}
&\text{Programming, proofs, probability}\longrightarrow\text{models (1)}
\longrightarrow\text{recurrences (2)}\\
&\longrightarrow\text{generating functions (3)}\longrightarrow\text{asymptotics (4)}
\longrightarrow\text{symbolic method (5)}.
\end{align*}
Chapters 6--9 depend on the methods in Chapters 2--5, but can be studied in application-driven order after Chapter 5. For a tree-focused learner, emphasize 2, 3, 5, 6; for search-system engineering, emphasize 1, 3, 4, 8, 9; for sorting analysis, emphasize 1, 2, 4, 7.
\section{Assessment coverage matrix}
\begin{longtable}{lll}\toprule Outcome & Form / artifact & Evidence \\ \midrule
Model validity & analysis-record-blank & Exact input law, unit and scope \\
Proof discipline & proof-assessment-blank & Base cases, invariant and quantifiers \\
Empirical design & experiment-plan-blank & Raw repetitions and reproducibility \\
Laboratory quality & lab-report-blank & Source revision, tested counter \\
Capstone readiness & capstone-review-blank & Technical/empirical trade-off \\\bottomrule\end{longtable}
\section{Source integrity}
These notes paraphrase, extend and teach the underlying mathematics in original form. The independent laboratories and forms are not official instructor materials and contain no scanned diagrams, chapter prose, or copied exercise solutions. Use the licensed book for the authoritative presentation and exercise numbering.

\end{document}
